\QHeading{第7题　用简便方法计算下面各题。}
\par\noindent


\CalcBlock[方法一：简便计算]{1}{\begin{aligned}[t]&\phantom{={}}\color{black}98\times 45\\&\color{answerred}=(100-2)\times45\\&\color{answerred}=100\times45-2\times45\\&\color{answerred}=4500-90\\&\color{answerred}=4410\end{aligned}}
\par
\MethodBlock{方法二：按数位拆分计算}{把一个因数拆成整十、整百数与余下部分之和。}{\begin{aligned}[t]&\phantom{={}}98\times 45\\
&=98\times (40+5)\\
&=98\times 40+98\times 5\\
&=3920+490\\
&=4410\end{aligned}}


\CalcBlock[方法一：简便计算]{2}{\begin{aligned}[t]&\phantom{={}}\color{black}125\times 72\\&\color{answerred}=125\times (8\times 9)\\&\color{answerred}=(125\times 8)\times 9\\&\color{answerred}=1000\times 9\\&\color{answerred}=9000\end{aligned}}
\MethodBlock{方法二：把特殊因数改成整百或整千}{先把25、125等改写成整百或整千除以一个数。}{\begin{aligned}[t]&\phantom{={}}125\times 72\\
&=(1000\div 8)\times 72\\
&=1000\times (72\div 8)\\
&=1000\times 9\\
&=9000\end{aligned}}

\MethodBlock{方法三：凑整后补偿}{把接近整十、整百的因数凑整，再补回或减去多算的部分。}{\begin{aligned}[t]&\phantom{={}}125\times 72\\
&=125\times (70+2)\\
&=125\times 70+125\times 2\\
&=8750+250\\
&=9000\end{aligned}}

\MethodBlock{方法四：按数位拆分计算}{把一个因数拆成整十、整百数与余下部分之和。}{\begin{aligned}[t]&\phantom{={}}125\times 72\\
&=72\times (100+25)\\
&=72\times 100+72\times 25\\
&=7200+1800\\
&=9000\end{aligned}}

\CalcBlock[方法一：简便计算]{3}{\begin{aligned}[t]&\phantom{={}}\color{black}38\times 29+38\\&\color{answerred}=38\times(29+1)\\&\color{answerred}=38\times30\\&\color{answerred}=1140\end{aligned}}
\MethodBlock{方法二：先处理括号或乘除，再凑整再减去多加的数}{}{\begin{aligned}[t]&\phantom{={}}38\times 29+38\\
&=1102+38\\
&=1102+40-2\\
&=1142-2\\
&=1140\end{aligned}}

\CalcBlock[方法一：简便计算]{4}{\begin{aligned}[t]&\phantom{={}}\color{black}56\times 125-56\times 25\\&\color{answerred}=56\times(125-25)\\&\color{answerred}=56\times100\\&\color{answerred}=5600\end{aligned}}
\MethodBlock{方法二：按运算顺序直接计算（对照）}{本法只作结果核对；本题要求简便计算，前面的简便解法仍保留。}{\begin{aligned}[t]&\phantom{={}}56\times 125-56\times 25\\
&=7000-1400\\
&=5600\end{aligned}}

\CalcBlock[方法一：简便计算]{5}{\begin{aligned}[t]&\phantom{={}}\color{black}856-(856-278)\\&\color{answerred}=856-856+278\\&\color{answerred}=0+278\\&\color{answerred}=278\end{aligned}}
\MethodBlock{方法二：先处理括号或乘除，再多减再加回}{}{\begin{aligned}[t]&\phantom{={}}856-(856-278)\\
&=856-578\\
&=856-580+2\\
&=276+2\\
&=278\end{aligned}}

\CalcBlock[方法一：简便计算]{6}{\begin{aligned}[t]&\phantom{={}}\color{black}784-(184-59)\\&\color{answerred}=784-184+59\\&\color{answerred}=600+59\\&\color{answerred}=659\end{aligned}}
\MethodBlock{方法二：按运算顺序直接计算（对照）}{本法只作结果核对；本题要求简便计算，前面的简便解法仍保留。}{\begin{aligned}[t]&\phantom{={}}784-(184-59)\\
&=784-125\\
&=659\end{aligned}}

\CalcBlock[方法一：简便计算]{7}{\begin{aligned}[t]&\phantom{={}}\color{black}57+75-57+75\\&\color{answerred}=(57-57)+(75+75)\\&\color{answerred}=0+150\\&\color{answerred}=150\end{aligned}}
\MethodBlock{方法二：交换、结合后先凑整}{保留每个数原来的加减号，先把能凑整的部分放在一起。}{\begin{aligned}[t]&\phantom{={}}57+75-57+75\\
&=(57-57)+75+75\\
&=0+75+75\\
&=75+75\\
&=150\end{aligned}}

\MethodBlock{方法三：先处理括号或乘除，再提取相同因数}{用乘法分配律反过来计算。}{\begin{aligned}[t]&\phantom{={}}57+75-57+75\\
&=132-57+75\\
&=75+75\\
&=75\times (1+1)\\
&=75\times 2\\
&=150\end{aligned}}

\CalcBlock[方法一：简便计算]{8}{\begin{aligned}[t]&\phantom{={}}\color{black}2400-402\\&\color{answerred}=2400-(400+2)\\&\color{answerred}=2400-400-2\\&\color{answerred}=2000-2\\&\color{answerred}=1998\end{aligned}}
\MethodBlock{方法二：补差法}{从减数先补到一个整数，再补到被减数，两段差相加。}{\begin{aligned}[t]&\phantom{={}}2400-402\\
&=(410-402)+(2400-410)\\
&=8+1990\\
&=1998\end{aligned}}

\CalcBlock[方法一：简便计算]{9}{\begin{aligned}[t]&\phantom{={}}\color{black}249+78-49+122\\&\color{answerred}=(249-49)+(78+122)\\&\color{answerred}=200+200\\&\color{answerred}=400\end{aligned}}
\MethodBlock{方法二：交换、结合后先凑整}{保留每个数原来的加减号，先把能凑整的部分放在一起。}{\begin{aligned}[t]&\phantom{={}}249+78-49+122\\
&=(78+122)+249-49\\
&=200+249-49\\
&=449-49\\
&=400\end{aligned}}

\MethodBlock{方法三：按运算顺序直接计算（对照）}{本法只作结果核对；本题要求简便计算，前面的简便解法仍保留。}{\begin{aligned}[t]&\phantom{={}}249+78-49+122\\
&=327-49+122\\
&=278+122\\
&=400\end{aligned}}

\par\noindent

\QHeading{第8题　解决问题。}

\Problem{① 一块长方形菜地，长68米，宽32米，给这块菜地四周围上篱笆，篱笆的总长度是多少米? 这块菜地的面积是多少平方米?}
\par


\MethodTitle{方法一}
\Answer{篱笆总长：}
\Answer{\Formula{\begin{aligned}[t]&\phantom{={}}(68+32)\times2\\&=100\times2\\&=200\text{（米）}\end{aligned}}}
\Answer{菜地面积：}
\Answer{\Formula{\begin{aligned}[t]&\phantom{={}}68\times32\\&=68\times(30+2)\\&=68\times30+68\times2\\&=2040+136\\&=2176\text{（平方米）}\end{aligned}}}
\Answer{答：篱笆的总长度是200米。}
\Answer{答：菜地的面积是2176平方米。}

\MethodTitle{方法二：周长分别算两条长和两条宽；面积用补长方形}
\Answer{周长不先合并长与宽；面积把长68米暂补成70米，再去掉多补的2米宽长条。}
\Answer{\textbf{篱笆总长度}\par \Formula{\begin{aligned}[t]&\phantom{={}}68\times 2+32\times 2\\
&=136+64\\
&=200\text{（米）}\end{aligned}}}
\Answer{\textbf{菜地面积}\par \Formula{\begin{aligned}[t]&\phantom{={}}68\times 32\\
&=(70-2)\times 32\\
&=70\times 32-2\times 32\\
&=2240-64\\
&=2176\text{（平方米）}\end{aligned}}}
\Answer{答：篱笆的总长度是200米。\par 答：这块菜地的面积是2176平方米。}
\Problem{② 文具店卖出钢笔和圆珠笔各40盒，每盒钢笔15支，每盒圆珠笔10支。卖出的钢笔比圆珠笔多多少支?}

\MethodTitle{方法一}
\Answer{\Formula{\begin{aligned}[t]&\phantom{={}}(15-10)\times40\\&=5\times40\\&=200\text{（支）}\end{aligned}}}
\Answer{答：钢笔比圆珠笔多200支。}

\MethodTitle{方法二：分别求两种笔的总支数}
\Answer{先算钢笔与圆珠笔各有多少支，再比较。}
\Answer{\Formula{\begin{aligned}[t]&\phantom{={}}15\times 40-10\times 40\\
&=600-400\\
&=200\text{（支）}\end{aligned}}}
\Answer{答：卖出的钢笔比圆珠笔多200支。}
